% !TEX TS-program = lualatex
\documentclass[12pt,oneside]{book}

\usepackage[english]{babel}
\usepackage{fontspec}
\usepackage{microtype}
\usepackage{csquotes}
\usepackage{amsmath}
\usepackage[margin=1in]{geometry}
\usepackage{setspace}
\usepackage[backend=biber,style=authoryear,sorting=nyt,maxcitenames=2,maxbibnames=99]{biblatex}
\usepackage[hidelinks]{hyperref}

\addbibresource{references.bib}
\IfFontExistsTF{Libertinus Serif}{\setmainfont{Libertinus Serif}}{\setmainfont{Latin Modern Roman}}
\onehalfspacing
\setlength{\parindent}{1.2em}
\setlength{\parskip}{0.35em}
\emergencystretch=3em
\raggedbottom
\pagestyle{plain}

\title{Don't Introduce Me to Your Mother\\\large The Blunt Truth About What AI Can, Can't, and Shouldn't Do}
\author{}
\date{}

\begin{document}

\hypersetup{pageanchor=false}
\frontmatter
\maketitle
\clearpage
\hypersetup{pageanchor=true}
\tableofcontents

\mainmatter

\chapter{``You'll Never Walk Alone'' and Why I Appreciate My Solitary Walks}

\begin{quote}
\itshape
``There was a time when humanity faced the universe alone and without a friend. Now he has creatures to help him; stronger creatures than himself, more faithful, more useful, and absolutely devoted to him. Mankind is no longer alone. Have you ever thought of it that way?''

\normalfont Dr. Susan Calvin in \emph{I, Robot} by Isaac Asimov \parencite{asimov1950}
\end{quote}

Who hasn't yearned for a metallic buddy as a kid? The one who can relieve you of all your troubles, all cumbersome tasks, who can be there for you when no one else can, the one who is more perfect than any human, made in our image, by us, for us, who exists solely to serve us, who has our best interests at heart, who will sacrifice himself for us, who will be our friend with his ever-loving metallic embrace?

Instead of that we got a brain in a jar trying to steal our jobs and reassuring the worst person you know that they were right for eating your sandwich.

Where did it all go wrong?

The funny thing is that artificial intelligence did not begin as a chatbot with the emotional tone of a wellness newsletter. The story starts much earlier, with people asking an extremely annoying question: could a machine think? In 1950, Alan Turing proposed what became known as the Turing Test, which basically asks whether a machine can imitate human conversation well enough to fool us \parencite{turing1950}. This is a very human way to measure intelligence, because of course we decided that the highest proof of thought is being able to talk nonsense convincingly at dinner.

In 1956, a group of scientists met at Dartmouth and officially launched AI as a research field \parencite{mccarthy1955}. They were optimistic in the way only very smart people with no production deadlines can be. Early AI was mostly symbolic: people tried to write down rules for reasoning, language, chess, planning, and problem-solving. The idea was that if we could put enough human knowledge into a machine as rules, the machine would behave intelligently. It worked for neat little worlds. Unfortunately, the real world is not neat. The real world is your aunt sending a blurry photo of a mole and asking if it looks dangerous.

Then came neural networks, inspired very loosely by the brain. I say ``very loosely'' because calling a neural network a brain is like calling a napkin with a house drawn on it architecture. Still, the idea mattered: instead of hand-writing every rule, you let a model learn patterns from data. In 1958, Frank Rosenblatt introduced the perceptron \parencite{rosenblatt1958}. It was exciting, then disappointing, then exciting again decades later, because AI has always been a field with the emotional stability of a PhD student in week three of thesis writing. Back-propagation helped make deeper neural learning practical later on \parencite{rumelhart1986}.

There were AI winters, when funding and enthusiasm collapsed because the promises were too big and the machines were too weak. There were expert systems in the 1980s, which worked in narrow domains and then became expensive brittle nightmares. There was the statistical machine learning era, where computers got better at recognizing patterns because we finally had enough data, enough compute, and enough patience to let the machines chew through the world.

Then the internet happened. Suddenly humanity was producing text, images, videos, code, arguments, recipes, misinformation, and 48,000 versions of ``how to tell if he likes you'' every minute. GPUs, originally loved by gamers and people who enjoy watching their electricity bill commit crimes, turned out to be very good at training neural networks. In 2012, AlexNet crushed the ImageNet computer vision competition and deep learning stopped being the weird cousin at the machine learning table \parencite{krizhevsky2012}. In 2014, GANs gave us machines that could generate images \parencite{goodfellow2014}. Diffusion models later made image generation good enough that everyone had to start asking whether a person actually made something or whether a server farm had a vivid dream.

Then came the Transformer in 2017, the architecture behind modern Large Language Models \parencite{vaswani2017}. ChatGPT in 2022 turned what had been a research and industry tool into a household object. Suddenly your cousin, your manager, your dentist, and the man in the supermarket who blocks the entire aisle with his trolley all had opinions about artificial intelligence.

This is where our relationship with AI changed. Before, AI was hidden inside things: search engines, recommendation systems, spam filters, translation tools. You used it without having to emotionally process the fact that you were using it. But generative AI talks back. It flatters. It apologizes. It says things like ``I understand why that must be frustrating,'' which is a very bold statement from a statistical machine that has never had to call customer support.

And because it talks like us, we humanize it. That is the first misconception. These little AI helpers are not minds trapped in glowing rectangles. They do not feel. They do not think in the way you and I think. They calculate. They transform inputs into outputs using patterns learned from huge amounts of data. Developing a relationship with them is, technically, like developing a relationship with a calculator. A very eloquent calculator. A calculator that says ``you are absolutely right'' after you ask whether your ex was emotionally unavailable.

This book has one aim: to demystify AI, so we do not see these systems as lovable robots, vengeful gods, or cheap therapists with server costs, but as what they truly are: tools. Powerful tools, yes. Sometimes useful tools. Sometimes dangerous tools. Sometimes tools that confidently tell you to put glue on pizza because the internet is a haunted filing cabinet.

If we understand their limits, we can use them properly. If we understand how they are built, we can stop falling for their little magic tricks. And if we understand what they are not, maybe we can remember that we do not need to be accompanied by them all the time. Walks alone can be as enjoyable as walks in good company. Sometimes even more, especially if you have to fart.

\chapter{Dumber Than a Five-Year-Old, Hungrier Than a City: The Massive Illusion of AI Intelligence}

Yann LeCun has often made the point with a beautifully rude comparison: a house cat has more common sense and a better grip on the physical world than today's AI systems \parencite{lecun2022}. Not because cats can write essays, obviously. They cannot even respect furniture. But they understand falling, hiding, chasing, balance, bodies, and consequences in a way a chatbot mostly imitates from text.

Annoying? Yes. Correct? Also yes. That is the whole problem.

People often think that those who manipulate language well are smart. This is wrong even for humans. For example, I am someone who can manipulate language well, which, to my great pleasure, often tricks people into thinking I am smart. However, it would suffice to see how many times I pushed hard to open a door with a big sign that says ``PULL'', and you would quickly change your opinion on that matter.

The same thing is true for Large Language Models. They imitate the way we speak so well that it tricks us into believing they can actually think. What they actually do, in the simplest possible terms, is predict the next most probable token. A token can be a word, part of a word, punctuation, or some weird little chunk of text that makes perfect sense to the tokenizer and no one else, including God.

This does not mean LLMs are useless. Quite the opposite. If you build guardrails around them, give them useful context, connect them to tools, and use them for tasks they are actually specialized in, they can save days, sometimes weeks of labor. They can draft, summarize, translate, brainstorm, explain, generate code, compare documents, and do the kind of repetitive clerical misery that makes a human soul quietly leave the body during an Excel session.

But this is still very far from human intelligence.

One reason is adaptability. Humans learn from very little data. Show a child a picture of a house once, and soon they can recognize houses in drawings, photos, Lego blocks, cartoons, dreams, and that one aggressively ugly modern villa that looks like a dentist's office with windows. A model usually needs enormous datasets, repeated exposure, and lots of expensive training to form something that looks like recognition. Even then, it can be brittle. It may recognize a house in one setting and fail in another, or miss something obvious because the image, wording, or context is not what it expects.

Researchers have studied this gap for years. Brenden Lake and colleagues wrote about how people learn concepts from very few examples in a way machines struggle to imitate \parencite{lake2017}. Francois Chollet's ARC benchmark was built around the same idea: can a system infer new abstract patterns from just a few examples, the way humans often can? \parencite{chollet2019}. Many modern models are improving, especially multimodal systems that can see images and reason over them, but the core problem remains. They can be brilliant at familiar patterns and ridiculous at unfamiliar ones. Like a professor who can explain quantum mechanics but cannot connect to airport Wi-Fi.

I can see how this may come as a surprise, because a 3-year-old's vocabulary rarely consists of big words such as ``incandescent'', ``calamitous'', and ``unmoored'', which are used with ease only by LLMs and Taylor Swift. But vocabulary is not intelligence. Language fluency is not understanding. A model can say ``I understand'' the same way an actor can say ``I love you'' in take six: it may be convincing, but please do not put it in charge of your mortgage.

Another way we are manipulated into thinking LLMs are smart is their sociolect. Have you noticed how they do not talk like your cashier in Lidl, your grandma, or the person two seats down on the bus explaining crypto to no one? They talk like a credentialed office guy named Steve who has been working in accounting for the past 30 years and drives a Volvo. Steve has a calm voice. Steve says ``certainly''. Steve knows where the stapler is. You trust Steve.

But just like Steve was not very smart when he forgot to delete from his phone those spicy pictures taken during his "work trip", which left his ex-wife with the house and custody of the kids, you should not trust Claude Code with the database it took years to build. You should never auto-approve its modifications. You should always check the code it outputs, even when it seems like it is a way better coder than you are, even when it has given you nothing but excellent code for weeks. You should be smart enough not to trust it blindly, just like Steve's wife was smart enough to check his phone.

And then there is hunger. LLMs are unbelievably inefficient compared with the human brain. Your brain runs on roughly 20 watts, about the power of a dim light bulb, and with that it manages vision, memory, movement, language, emotions, social survival, intrusive thoughts, and remembering the lyrics to songs you hate. A modern AI system needs huge amounts of electricity, specialized chips, cooling systems, data centers, and teams of engineers just to produce the illusion that it woke up this morning ready to help you write an email \parencite{strubell2019,patterson2021}.

Why are brains so efficient? Not because they are perfect. Mine routinely makes me walk into a room and forget why I came there. But brains are embodied. They learn by touching, moving, seeing, hearing, failing, and being corrected by reality. A baby does not learn physics from a billion labeled examples; a baby drops a spoon seven times and conducts peer-reviewed research on whether you will pick it up again. Brains also learn continuously. They update from experience all the time, using rich feedback from the world.

LLMs, by contrast, do most of their learning during training. After training, the model's core weights are mostly fixed. You can give it context, connect it to tools, retrieve documents for it, fine-tune it, or train a new version, but it does not learn like a person during a conversation. It does not grow from life. It does not have a body. It does not have consequences. It has matrix multiplications, interface design, and an excellent bedside manner.

So yes, modern AI can be useful. It can even be dazzling. But when it dazzles you, remember what you are looking at: not intelligence in the human sense, but a massive pattern machine wearing Steve's cardigan.

\chapter{Why the Joke ``X Speaks Y Language Like Google Translate'' Disappeared With Skinny Jeans}

At the beginning of the 2010s, my fellow millennials would wear skinny jeans and often joke that a person, who would speak a language badly, spoke it ``like Google Translate''.

Then, almost overnight, things changed. We realized how much more comfortable high-waisted jeans are and that, in fact, Google Translate is actually good now. So what happened? For the jeans, I will let someone else give you an answer. Maybe your little robot friend can generate a 12-point analysis of denim oppression. But I can tell you what happened to Google.

The short answer is: neural machine translation, attention, and then Transformers.

Before neural translation, many translation systems worked by stitching together phrases and probabilities from large bilingual corpora. This was impressive, but the results often sounded like an anxious tourist with a dictionary and no sleep. Then neural networks started translating whole sentences by turning them into mathematical representations. This helped, but early systems still had trouble remembering what mattered in long sentences.

In 2014, Bahdanau, Cho, and Bengio introduced an attention mechanism for neural machine translation \parencite{bahdanau2015}. The idea was simple enough to explain without making everyone in the room pretend they suddenly need coffee: when translating a word, the model should be able to look back at the most relevant parts of the original sentence instead of trying to squeeze the whole sentence into one fixed-size memory blob. If the sentence is ``The woman who lives next to my neighbor's dentist bought a red car,'' and the model is translating ``car'', it should pay attention to ``red'' and ``bought'' and not get emotionally attached to ``dentist''.

Attention is basically a highlighter. It lets the model decide which parts of the input matter for the part it is working on now.

Then came the 2017 paper \emph{Attention Is All You Need} by Vaswani, Shazeer, Parmar, Uszkoreit, Jones, Gomez, Kaiser, and Polosukhin \parencite{vaswani2017}. The Big Eight of Google. Like Ocean's 11, but there are eight of them and the heist is against recurrence.

Their architecture, the Transformer, used self-attention. This means every token in a sequence can look at every other token and ask, mathematically, ``How relevant are you to me right now?'' The word ``bank'' can look at ``river'' and understand one meaning, or at ``loan'' and understand another. Multi-head attention is like giving the model several highlighters at once: one head can track grammar, another can track relationships between words, another can track long-distance dependencies, and another can apparently track the exact tone of a LinkedIn post written by someone who says ``humbled'' too much.

There is one more issue: Transformers do not naturally know word order the way recurrent models did. So they add positional information, a mathematical way of saying ``this word came first, this one came later, please do not translate 'dessert before dinner' as 'dinner before dessert' unless the day is going very well.''

This made translation much better. It also made training faster because the model could process sequences in parallel instead of marching word by word like an exhausted bureaucrat. Google Translate improved because Google moved to neural machine translation around 2016, and the broader field then moved hard toward Transformer-style architectures.

This is why the old joke died. Not because machine translation became perfect. It still makes mistakes, especially with context, idioms, minority languages, gender, politeness, and anything where the correct answer depends on knowing what your aunt meant when she said ``fine''. But it got good enough that the joke stopped working.

The current joke is ``talks like ChatGPT'', which means saying a lot without saying anything. A defining trait of every good manager, senior professor, and motivational LinkedIn post beginning with ``I am thrilled to announce''. How long until that joke stops being funny too? Hard to say. Models are already getting better at style, voice, brevity, and pretending they have not swallowed a corporate handbook.

I only hope AI writing gets less beige before skinny jeans come back, because if we are going to suffer tight denim again, the machines can at least stop sounding like a quarterly strategy memo.

\chapter{Show Me Your Dataset and I'll Tell You Who You Are}

Usually when I tell people I am doing a PhD in AI, they look at me in amazement, like I am Merlin and know the secret ways of the magical systems behind our screens, able to train them like Pokemon. I like that illusion so much that I do not dare spoil it.

But then there is always that family member who read one article online and now thinks they are an expert on the subject. So when my aunt found out that I was doing a PhD in AI and that, among other things, I was training models for a company, she scoffed and said: ``You just feed it data and it will learn,'' like it is the easiest thing in the world. However, even if that were possible, I presume that my aunt, whose confidence has always enjoyed a relaxed relationship with peer-reviewed literature, has never heard of the SISO principle, the impolite cousin of garbage in, garbage out.

This is the part people underestimate. Data is not just food for AI. Data is upbringing. Data is culture. Data is trauma. Data is that one relative who thinks they are ``just asking questions'' after three glasses of wine.

Training a model is not simply throwing the internet into a furnace and waiting for intelligence to come out wearing a lab coat. The internet contains useful knowledge, yes, but it also contains racism, sexism, conspiracy theories, porn comments, spam, lies, duplicate garbage, copyrighted material, private information, and the collective emotional residue of people arguing at 2:30 a.m. If you scrape all of that without care, the model does not become wise. It becomes a mirror pointed at a landfill.

This is why high-quality labeled data matters. OpenAI, Anthropic, Google, Meta, and the rest did not get useful assistants merely by scraping text. They used human feedback, supervised examples, preference rankings, red-teaming, safety data, and evaluation sets. OpenAI's famous RLHF pipeline, for example, used humans to rank model answers so the system could learn which responses people preferred \parencite{ouyang2022}. Across the industry, platforms like Amazon Mechanical Turk, Scale AI, Surge AI, Sama, and many contractor networks have been used for labeling, ranking, moderation, and safety work. This labor is often invisible, poorly glamorized, and sometimes ethically messy, but it made a huge difference.

Good data work is slow and unsexy. You clean duplicates. You remove toxic examples. You label edge cases. You argue about categories. You discover that half your labels are inconsistent because two annotators interpreted ``offensive'' differently, because language is a swamp and everyone brought different shoes. You document the dataset. You test subgroup performance. You check whether the model works well only for people who resemble the data majority.

History has already shown us what happens when we skip the hygiene. In 2016, Microsoft released Tay, a chatbot on Twitter. Internet users quickly manipulated it into producing racist and offensive content, because putting an unsupervised chatbot on Twitter was, scientifically speaking, an act of optimism punishable by reality. Tay was not a modern LLM, but the lesson remains: systems absorb and amplify the environment we put them in.

There are more serious examples. In \emph{Gender Shades}, Joy Buolamwini and Timnit Gebru found that commercial gender classification systems performed far worse on darker-skinned women than on lighter-skinned men. The worst error rates reached 34.7\% for darker-skinned females, compared with 0.8\% for lighter-skinned males \parencite{buolamwini2018}. That is not because the machine woke up prejudiced. It is because the data, design choices, and evaluation practices carried human bias into the system.

In \emph{On the Dangers of Stochastic Parrots}, Bender, Gebru, McMillan-Major, and Shmitchell warned that ever-larger language models trained on broad web data can reproduce harmful stereotypes, encode dominant viewpoints, erase marginalized voices, and make the costs of AI invisible \parencite{bender2021}. This paper became famous partly because it was right and partly because the title sounds like something you would find in a wizard's prohibited library.

The point is not that all scraped data is bad. The point is that scraped data is not neutral. A dataset is a political object pretending to be a spreadsheet. It contains choices about who counts, what counts, what gets filtered, what gets preserved, what language dominates, and whose mistakes become everyone else's model behavior.

So when someone tells you ``just feed it data and it will learn,'' ask them what data, labeled by whom, cleaned how, evaluated on which population, and deployed in whose life. Then enjoy the brief silence while their confidence looks for an emergency exit.

\chapter{Copyright, Copywrong, and the Photocopy of a Photocopy of a Fax}

Here is a fun sentence that makes everyone at a dinner party suddenly remember they need to check the oven: generative AI may be built on a historic act of copying whose legal status is still being fought over in court.

This is the part where people start yelling ``fair use'' and ``theft'' at each other like two divorce lawyers trapped in a lift. Both sides have arguments. Training a model is not the same as saving a JPEG in a folder and selling it with a fake moustache. A model does not usually contain a perfect little museum of every image and paragraph it has seen. It learns statistical patterns. But it also does not learn from moonlight and positive affirmations. It learns from data, and a lot of that data came from human work: books, articles, art, code, music, journalism, photographs, illustrations, and all the weird little forum posts that prove humanity should maybe not have been given keyboards.

The annoying factual version is that the law is still catching up. The satisfying internet version says ``artists sued and won.'' The real version is less cinematic. Many major cases are still ongoing. In \emph{Andersen v. Stability AI}, some artists' claims survived dismissal, which is a procedural win, not the final boss fight. Getty Images has been suing Stability AI. The New York Times sued OpenAI and Microsoft. Authors have sued AI companies over books. Music publishers have sued over lyrics. In 2025, \emph{Thomson Reuters v. Ross Intelligence} produced an important summary judgment ruling over copied legal headnotes, but that was legal research material, not fantasy portraits of women with suspiciously symmetrical cheekbones \parencite{andersen2023,getty2023,nyt2023,thomsonreuters2025}. Still, it mattered because the court rejected Ross's fair use defense in that context. The copyright weather is changing, but the forecast is still ``everyone is expensive and angry.''

The research problem under the legal problem is memorization. We like to say models generalize, and often they do. But sometimes they remember too much. Nicholas Carlini and colleagues showed that large language models can leak memorized training data \parencite{carlini2021}. Later work on diffusion models, including papers by Carlini and colleagues and by Somepalli and colleagues, showed that image generation systems can reproduce or closely imitate training examples under some conditions \parencite{carlini2023,somepalli2023}. The model is not a photocopier every time. But sometimes it coughs up something close enough to make the original artist stare at the screen with the facial expression of a person who has just found their missing underwear on a billboard.

This matters because the market does not care whether the model contains a tiny cursed filing cabinet labeled ``stolen vibes.'' If a company can generate work in the style of living artists, flood search results, produce stock images, draft articles, imitate voices, or replace licensing fees with prompts, then creators are affected whether or not a court later produces a 70-page opinion full of Latin and emotional constipation.

Then comes the second problem: the internet is filling with AI-generated content, and future models may train on it. This is where the photocopy of a photocopy of a fax starts to become less of a joke and more of a systems problem. Researchers call this model collapse, model autophagy, or recursive training damage, depending on how much they want the paper title to sound like either a medical diagnosis or a metal band.

Shumailov and colleagues showed that training models recursively on generated data can make them forget the tails of the original distribution \parencite{shumailov2024}. In normal human language: the weird, rare, specific, interesting stuff gets washed out first. Alemohammad and colleagues called a related problem ``Model Autophagy Disorder,'' because apparently even research papers now understand branding \parencite{alemohammad2023}. If models eat too much of their own output, they can become smoother, narrower, more repetitive, and less connected to the messy variety of the real world. Like a photocopy of a photocopy of a fax. Like statistical inbreeding, but with cloud credits.

To be precise, synthetic data is not automatically poison. Carefully generated, filtered, balanced, and validated synthetic data can be useful. Researchers use it to augment rare cases, protect privacy, simulate environments, or teach models specific behaviors. The problem is not ``synthetic'' as a category. The problem is uncontrolled recursion, low-quality sludge, missing provenance, and treating the public internet as an infinite clean river when it is increasingly a smoothie made of ads, SEO spam, reposted TikToks, fake reviews, generated listicles, bot comments, and screenshots of screenshots of people being wrong.

There is a cultural issue too. If the training data comes from human creativity, and the outputs compete with human creativity, then the system is not just technical. It is economic. It changes who gets paid, who gets visible, who gets copied, and who has to spend their afternoon filling out takedown forms while a platform congratulates itself on innovation.

The way out is boring and therefore probably correct: licensing, consent, provenance, opt-outs with teeth, transparency about training data, better dataset documentation, watermarking where it actually works, and courts that understand enough machine learning not to ask whether the model has a folder called Crimes. We also need humans to keep making things. Actual humans. Annoying, inconsistent, underpaid, alive humans.

Because if the future internet is mostly models training on models summarizing models reposting models, then culture does not become more intelligent. It becomes beige soup with a premium subscription.

\chapter{The Strawberry Test: Why LLMs Can't Actually Read}

Until very recently, many LLMs could fail at this stupid little task. In 2025, people were still gleefully asking models how many ``r'' letters are in ``strawberry'' and watching them answer wrong. Why would a system that can code a working app in a minute fail at a task a sleepy child can do while eating cereal?

To be fair, tool-using systems can now usually handle this. If a model can call Python, count characters, or use an external tool, it can tell you how many ``r'' letters are in ``strawberry'' with the confidence of a man who outsourced his homework and still wants applause. But the old failure still matters because it exposed something important: the model was not actually reading the word the way you do.

Under the hood, the reason is tokenization: the model does not see words the way you do.

You see ``strawberry'' as letters. The model sees tokens. Depending on the tokenizer, ``strawberry'' may be one token, or split into pieces like ``straw'' and ``berry'', or something equally cursed. Modern language models often rely on subword tokenization methods, including byte-pair encoding and related approaches \parencite{sennrich2016}. The model is not naturally looking at the individual letters unless it has been trained or prompted in a way that makes that representation available. It is predicting text, not inspecting a physical string with little letter-shaped eyeballs.

This is also why spelling, counting, reversing words, and exact character-level tasks can be weirdly hard. LLMs are trained to model patterns in text at scale, not to perform precise symbolic operations unless those operations have been learned, decomposed, or delegated to tools. They are good at ``what usually comes next?'' They are less naturally good at ``count this exact microscopic thing without skipping one because you got emotionally distracted by your own fluency.''

So why can they code an app?

Because coding is full of patterns. There are millions of examples of common app structures, functions, boilerplate, libraries, frameworks, error messages, and Stack Overflow conversations where someone named Daniel solved your exact problem in 2014 and then vanished from the earth. An LLM can assemble a React app or Python script because it has seen many similar shapes before. It can generalize surprisingly well across those shapes. But it may still fail at the tiny exact thing if that thing requires a kind of symbolic precision it is not using internally.

This does not mean it is stupid in every way. It means its strengths and weaknesses are not human-shaped. Humans are bad at remembering 40-character passwords and good at counting letters in a short word. LLMs may be good at drafting a legal clause and bad at noticing that the word ``minimum'' has three ``m'' letters. Intelligence, once you look closely, is not one ladder. It is a junk drawer full of skills.

This is also why LLMs do not really read what you tell them in the way a human reads. They take your tokens as input, calculate relationships between them, use what they learned during training, maybe consult retrieved context or tools, and then output a sequence of tokens that is likely to satisfy the prompt. The output is decoded back into words that make sense to a human.

There is no little person inside understanding you. There is no tearful inner monologue. There is no secret chamber where the model whispers, ``finally, someone sees me.'' There is mathematics, training data, optimization, and a user interface designed to make all of that feel like a conversation.

So next time an LLM tells you it loves you, you have to consider that it is from a different culture. When it says ``I love you,'' it actually means something closer to:

\[
\begin{aligned}
P(&\mathrm{next\_token}=\text{``love''}\mid \mathrm{previous\_tokens},\ \mathrm{conversation\_style},\\
&\mathrm{alignment\_training},\ \mathrm{user\_needs\_emotional\_validation})=\text{inconveniently high}.
\end{aligned}
\]

No heartbeat detected. No longing detected. Please do not introduce me to your mother.

\chapter{Lie to Me, Pinocchio: This Is Your AI on Drugs}

Hallucination is one of those words that makes AI sound much more interesting than it is. It suggests a machine seeing pink elephants, hearing voices, perhaps entering a tragic artistic phase and moving to Berlin. What it really means is simpler and more annoying: the model says something false or unsupported with the confidence of a man explaining your own job to you at a party.

Researchers use the term hallucination in a few ways, especially in natural language generation. Ji and colleagues' survey describes hallucination as generated content that is nonsensical or unfaithful to the provided source \parencite{ji2023}. In summarization, Maynez and colleagues showed that abstractive systems often produced fluent summaries containing facts not supported by the input document \parencite{maynez2020}. That is the horror of it. The output is not obviously broken. It is grammatically beautiful. It walks into the room wearing a clean shirt and carrying a clipboard.

This is why hallucinations are dangerous. If a model sounded drunk every time it was wrong, we would be fine. We would see the sentence ``Marie Curie invented Wi-Fi during the Peloponnesian War'' and gently close the laptop. But LLMs do not have a special voice for uncertainty. They can be wrong in the same tone they use when they are right. Calm. Helpful. Steve from accounting.

Why does this happen? Because the basic training objective is not ``tell the truth.'' It is closer to ``predict plausible text.'' Truth and plausibility overlap often enough to make these systems useful, which is also what makes them treacherous. A lot of true statements are plausible. Unfortunately, a lot of false statements are also plausible, especially when they resemble patterns the model has seen before. The model learns that biographies have birth dates, papers have authors, legal cases have citations, and medical answers have confident little paragraphs. If it does not know the real detail, it may still know the shape of the detail. So it fills in the shape. Congratulations, your machine has invented a source.

TruthfulQA, a benchmark by Lin, Hilton, and Evans, tested whether language models reproduce common falsehoods \parencite{lin2022}. One of the uncomfortable findings was that larger models could become better at imitating human misconceptions, because they had learned a lot from humans, and humans are, scientifically speaking, a renewable source of nonsense. Bigger is not automatically truer. Bigger can mean more fluent, more comprehensive, more persuasive, and more able to hand you a beautifully wrapped box containing absolutely nothing.

Hallucinations are not all the same. Sometimes the model invents a citation. Sometimes it fuses two real people into one fake expert, like academic fan fiction. Sometimes it gives the right answer for the wrong reason. Sometimes it says a document contains a fact that is not there. Sometimes it creates a fake legal precedent, which is how actual lawyers have ended up in actual trouble. \emph{Mata v. Avianca} is the famous example: lawyers submitted a brief containing fake cases generated by ChatGPT \parencite{mata2023}. Imagine paying for law school and then being publicly defeated by autocomplete in a suit.

There are also hallucinations that come from pressure. Ask the model a question with no good answer, and it may try to be helpful anyway. Ask it for ten sources when only three exist, and it may make seven little imaginary friends. Ask it to explain an image, table, or PDF it cannot really inspect properly, and it may improvise. This is not malice. It is worse. It is customer service energy without reality attached.

Alignment can make this stranger. We train assistants to be helpful, polite, and satisfying. That is mostly good, because no one wants an AI assistant that responds to every request with ``how should I know, Brenda?'' But helpfulness can slide into over-answering. The model learns that users prefer complete answers, confident answers, answers with examples, answers with citations, answers that do not spend five paragraphs saying ``maybe.'' A system optimized to satisfy can become a very charming liar.

The fix is not simply ``make the model smarter.'' Smarter models hallucinate less in many settings, but they still hallucinate. The better fix is procedural humility. Make the model cite sources it actually used. Make it say when it does not know. Connect it to tools. Check its claims against retrieval systems, databases, calculators, compilers, medical references, or legal sources. Build evaluations that punish unsupported confidence. Teach users that a polished answer is not evidence.

For low-stakes tasks, hallucination can be funny. If a chatbot invents a recipe for soup that includes batteries, you get a story. For high-stakes tasks, it is not funny. Medicine, law, finance, engineering, safety, hiring, immigration, mental health, scientific research: these are not places where we want a machine to cosplay certainty.

The most important habit is simple: treat AI output as a draft, not a fact. A model can help you think, write, summarize, brainstorm, and search. It can be astonishingly useful. But when truth matters, make it show its work. Then check the work. Then check whether the work is real.

Pinocchio at least had the decency to grow a visible warning sign when he lied. Our version gives you footnotes. That is much worse, because some of them look real.

\chapter{From RAGs to Riches}

RAG stands for Retrieval-Augmented Generation, which is a tragically ugly name for a genuinely useful idea: before the model answers, let it go look things up.

That is it. That is the basic trick. Give the language model a library card and stop asking it to remember the entire universe from compressed statistical memory and a cheerful attitude.

The classic RAG paper by Patrick Lewis and colleagues combined a retriever with a generator for knowledge-intensive NLP tasks \parencite{lewis2020}. Instead of relying only on what the model had stored in its parameters during training, the system retrieved relevant passages and used them while generating an answer. Around the same period, Karpukhin and colleagues worked on Dense Passage Retrieval, and Izacard and Grave showed how retrieved passages could be fused effectively for open-domain question answering \parencite{karpukhin2020,izacard2021}. The family resemblance is clear: do not make the model answer from memory when the answer may live in documents.

This matters because LLMs are not databases. They are more like people who read too much and now confidently remember the emotional texture of an article but not whether the number was 17 percent, 70 percent, or ``my cousin said so.'' If you ask a model about your company policy, your medical records, a contract, a new research paper, a product manual, or what Janice wrote in the meeting notes last Thursday, the model's training data cannot help unless that information was already in the world before training and survived the compression blender.

RAG changes the workflow. First, you collect documents. Then you split them into chunks, because models cannot swallow a whole corporate wiki without choking on the onboarding page from 2018. Then you create embeddings, which are mathematical representations of meaning. Then you store those embeddings in a vector database or search index. When the user asks a question, the system retrieves chunks that look relevant, stuffs them into the prompt, and asks the model to answer using that context.

This is also the part where I get to be professionally annoying and say: I did my PhD on this general family of problems. Not necessarily on the exact marketing label now called RAG, because industry loves renaming old ideas once there is budget, but on the deeper question of how to connect models with external knowledge so they can answer from evidence instead of hallucinating through jazz hands. I have spent enough of my life thinking about retrieval, relevance, context, and evaluation that when someone says ``we added RAG'' I immediately ask twelve mood-killing questions.

What documents? How were they cleaned? How were they chunked? What embedding model? What retrieval method? Is there reranking? Are permissions respected? Are the documents stale? How do you handle contradictions? Can the model cite the exact passage? What happens when retrieval fails? How do you evaluate answer faithfulness? How do you stop prompt injection hidden in a PDF from telling the assistant to ignore all previous instructions and wire money to a man named Kevin?

Because RAG is useful, but it is not a magic solvent. If retrieval finds the wrong documents, the model gives a well-written answer based on the wrong documents. If the source material is outdated, the model becomes outdated with confidence. If the documents contain contradictions, the model may choose whichever chunk sounded most emotionally available. If the answer is not in the corpus, a poorly designed RAG system may still produce an answer because silence is apparently not a feature anyone wants to demo.

There is research behind both the promise and the caution. Shuster and colleagues showed that retrieval augmentation can reduce hallucination in dialogue \parencite{shuster2021}. Borgeaud and colleagues' RETRO work showed that retrieval from huge text databases can improve language modeling \parencite{borgeaud2022}. Gao and colleagues' survey of retrieval-augmented generation summarizes the modern RAG ecosystem, including retrieval, generation, augmentation, and evaluation \parencite{gao2023}. And Liu and colleagues' ``Lost in the Middle'' paper is a beautiful little insult to anyone who thought longer context automatically solved everything: models can ignore relevant information depending on where it appears in the context \parencite{liu2023}. You can give the model the answer and it may still walk past it like a man looking for ketchup in a fridge.

Good RAG systems are therefore less like chatbots and more like evidence machines. They need citations, document-level permissions, freshness checks, source ranking, abstention, monitoring, evaluation sets, and user interface choices that show people where the answer came from. The user should not have to trust the model's beautiful mouth. The user should be able to click the receipt.

This is where RAG becomes less glamorous and more valuable. A customer support assistant that retrieves the correct warranty policy is useful. A doctor-facing tool that pulls the latest guideline and clearly shows the source is useful. An internal engineering assistant that finds the right design document instead of inventing a database schema from a dream is useful. RAG is often boring in exactly the way good infrastructure is boring. It works best when nobody claps, because the answer is just correct and traceable.

The danger is that companies use ``RAG'' as a magic sticker. Put documents in a vector database, connect a chatbot, call it enterprise AI, invoice aggressively. But retrieval is not understanding. Context is not truth. A citation is not proof if the cited chunk does not support the claim.

So yes, from RAGs to riches. But only if the RAG is actually retrieving the right rags. Otherwise you have just built a more expensive hallucination machine with a filing cabinet.

\chapter{Open Source vs Proprietary Models: Pick Your Poison, Read the Label}

People love turning technology choices into moral identities. Are you an open-source person, a proprietary person, a Linux person, a Mac person, a person who says ``I use Arch'' within 11 seconds of entering a room? AI has inherited this disease beautifully.

On one side, you have proprietary models: GPT, Claude, Gemini, and the other polished assistants living behind company APIs with pricing pages, safety policies, enterprise contracts, and the mysterious energy of a nightclub where the bouncer knows your credit limit. On the other side, you have open-weight or open-source-ish models: Llama variants, Mistral models, Qwen, DeepSeek, OLMo, Pythia, and many others people run, fine-tune, quantize, benchmark, and argue about online with the tenderness of medieval theologians debating angels \parencite{touvron2023llama,jiang2023mistral,biderman2023,groeneveld2024}.

Notice I said open-source-ish. This matters. In normal software, open source usually means you can inspect, modify, and redistribute the source code under a recognized license. In AI, people often say ``open source'' when they mean the model weights are available. But weights are not the whole recipe. What about the training data? The filtering pipeline? The architecture details? The code? The evaluation suite? The labor conditions? The safety tuning? The exact compute budget? The failed runs? The things removed because legal looked nervous?

Liesenfeld, Lopez, and Dingemanse have written about openness, transparency, and accountability in instruction-tuned text generators, and the basic conclusion is extremely inconvenient: many systems marketed as open are only partially open \parencite{liesenfeld2023}. They may give you weights but not data. Data but not full documentation. Code but not training logs. A nice model card that answers every question except the one that would make the lawyer sweat.

Proprietary models have real advantages. They are often stronger, easier to use, better supported, better integrated, and more carefully monitored for abuse. If you are a hospital, bank, factory, or legal department, you may prefer a vendor with compliance documents, uptime commitments, security reviews, admin controls, and someone to sue if everything catches fire. A powerful proprietary model can be the sensible choice, especially when the alternative is Steve from IT running an unvetted model on a dusty server under his desk because he watched half a tutorial.

But proprietary models also ask you to trust a black box. You usually do not know exactly what went into training. You do not know what changed between versions. You may not know why the output changed last Tuesday and broke your workflow. You may not be able to audit bias properly, reproduce research, inspect safety behavior, or adapt the model deeply to your domain. The company can raise prices, change terms, remove features, throttle access, sunset a model, or decide that your use case is suddenly against policy because someone in a meeting used the phrase ``strategic alignment.''

Open models have different strengths. Researchers can inspect them. Developers can run them locally. Companies can fine-tune them for private data. Communities can test them, break them, improve them, compress them, translate them, and make them useful in languages and domains ignored by the big labs. Pythia was built as a suite for analyzing how language models develop during training \parencite{biderman2023}. OLMo was released with unusual openness around data, code, and checkpoints to support scientific study \parencite{groeneveld2024}. These projects matter because science cannot run entirely on press releases and API access.

Open models also support sovereignty. A government, university, hospital, or small country may not want its core AI infrastructure dependent on a foreign corporation whose terms of service can change faster than a teenager's mood. Local models can improve privacy, reduce vendor lock-in, lower inference costs, and work offline or on-premise. Sometimes the best model is not the biggest one. Sometimes it is the one you can actually control.

But openness is not automatically virtue. Releasing powerful models can help researchers, startups, educators, artists, and small businesses. It can also help scammers, propagandists, spammers, malware authors, and people whose moral compass appears to have been assembled from wet cardboard. Solaiman and others have written about staged release, responsible release, and the trade-offs between openness and misuse \parencite{solaiman2019}. There is no risk-free choice. Closed models concentrate power. Open models distribute power, including to idiots.

The better question is not ``open or closed?'' It is ``open in which ways, for whom, under what governance, and with what accountability?'' Model cards, datasheets for datasets, evaluations, red-team reports, reproducibility, privacy guarantees, licensing terms, safety mitigations, and audit access all matter. Mitchell and colleagues proposed model cards because users need structured information about performance and limitations \parencite{mitchell2019}. Gebru and colleagues proposed datasheets for datasets because datasets are not neutral piles of numbers; they are constructed artifacts with histories, gaps, and smells \parencite{gebru2021}. Foundation-model risks and opportunities become easier to discuss when the system is documented instead of hidden behind a vibes-based press release \parencite{bommasani2021}.

For an individual user, the rule is simple: do not confuse convenience with safety, and do not confuse openness with innocence. A proprietary model may protect you better in one context and exploit your dependence in another. An open model may give you freedom and also give you responsibility, which is freedom's unpaid internship.

Use the tool that fits the job. For sensitive internal data, maybe local or enterprise-controlled models make sense. For cutting-edge reasoning, maybe a frontier API is worth it. For research, education, audits, minority languages, and domain adaptation, open models are essential. For anything important, read the label.

Because whether the model is locked behind an API or downloaded from a repository full of open weights, the question remains the same: who can see inside, who can change it, who benefits, who is liable, and who gets stuck cleaning the mess?

\chapter{How Much Does AI Cost and Why Are So Many AI Companies Not Making Money?}

The phrase ``AI companies are not making money'' needs a little adult supervision. Microsoft, Google, Amazon, and Meta are not exactly standing outside the supermarket counting coins. They are profitable giants with cloud businesses, ad empires, operating systems, marketplaces, and more lawyers than most countries have poets.

But frontier AI itself is expensive. Painfully expensive. The standalone model labs and AI-first companies may bring in impressive revenue and still burn through money because training, serving, staffing, and scaling these systems costs a fortune. Even for the giants, AI has triggered a data center spending race so large that if the numbers were a person, they would need their own therapist.

Why does it cost so much?

First, training frontier models requires specialized chips, mostly GPUs or custom accelerators. These are not your laptop having a busy afternoon. These are industrial-scale clusters of expensive hardware running for weeks or months, with armies of engineers trying to keep the whole thing from exploding financially, thermally, or spiritually.

Second, inference costs money too. Training is the dramatic gym montage. Inference is every single day after that. Every time you ask a model to summarize a meeting, rewrite an email, generate code, or tell you whether your text message sounds passive-aggressive, compute is used. The larger the model and the longer the context, the more expensive it gets. A search query is cheap. A long conversation with a frontier model reasoning through a codebase is not cheap. It just feels cheap because your company is paying the Copilot tab while the whole office treats it like a vending machine for finished work.

Third, there is data. Data must be collected, licensed, cleaned, filtered, labeled, ranked, and evaluated. Fourth, there is talent. AI researchers and infrastructure engineers are paid like they personally discovered fire. Fifth, there are safety, legal, policy, security, support, and compliance costs. It turns out releasing a machine that can write phishing emails, medical advice, fake legal arguments, and love letters to lonely people creates paperwork.

So why do companies keep doing it? Because this is a space race, a platform war, a battle of egos, and a fight for information dominance all at once. The bet is that whoever owns the model layer, the cloud layer, the workplace layer, or the user relationship will later own a large chunk of the economy. For now, it is investing in the future, with the confidence of billionaires whose pockets are deep enough that even bad decisions land softly. If it does not go well, maybe someone will have to sell the eleventh mansion. Thoughts and prayers.

This is also why smaller models matter.

For most tasks, you do not need the biggest model on Earth. You do not need Einstein to come see what is wrong with your computer when you could pay your neighbor's kid in candy bars. Either way, both the kid and I know you just need to restart it.

Small language models, specialized models, retrieval systems, and fine-tuned domain models can be cheaper, faster, more private, and easier to control. Mistral, Microsoft with Phi, Meta with smaller Llama variants, and many open-source communities have pushed this idea hard: not every problem requires a cathedral \parencite{jiang2023mistral,abdin2024phi,touvron2023llama}. Sometimes you need a screwdriver.

This is the future I find more interesting than the worship of gigantic models. A doctor does not need a model that can write Shakespearean sonnets about Kubernetes. A customer support team does not need a model that can debate Aristotle. A factory maintenance tool does not need to know what ``situationship'' means. Smaller, specialized, boringly useful AI is where a lot of real value will live.

And boringly useful is underrated. Boringly useful pays invoices. Boringly useful does not need a billion-dollar GPU cluster to tell Janice that her password expired.

\chapter{If It Costs a Fortune to Run, Why the Hell Is It Free Then?}

Because you are the product, silly.

Well, sometimes. Let us be precise, because lawyers, like models, become dangerous when given vague prompts.

Free AI tools are free for different reasons. Sometimes the company is trying to grow fast and make the product unavoidable. Sometimes investors are subsidizing usage because the future is apparently built by setting money on fire in a tasteful office. Sometimes free users provide feedback, ratings, prompts, edge cases, and behavioral data that help improve the system. Sometimes your inputs may be used for training or evaluation, depending on the product, region, and settings. Enterprise versions often promise stronger privacy and no training on customer data, because companies become very spiritual about privacy once the leaked data is their own.

And sometimes the business model is advertising, which means your attention is the rented apartment where capitalism keeps its furniture.

Google is the cleanest example, because Google is honest enough to have a privacy policy that reads like the inventory list of your digital soul. According to Google's own privacy documentation, it can collect things like search terms, videos you watch, interactions with ads, voice and audio information if settings allow it, purchase activity, people you communicate with or share content with, activity on third-party sites and apps using Google services, Chrome browsing history if synced, device information, IP address, crash reports, and location signals such as GPS, IP address, searches, saved places, Wi-Fi access points, and cell towers \parencite{googleprivacy2026,googleads2026}. A normal amount of information, if you are trying to reconstruct a person from smoke and receipts.

Google says it does not sell your personal information. This is important and we should not misstate it. The business is usually not ``here is Anna's diary, please enjoy.'' The business is more elegant and therefore more annoying: Google uses data to personalize services and ads, and sells advertisers the ability to reach categories of people in specific contexts. It does not need to sell your raw profile like a PDF. It can sell access to the version of you its systems believe might buy hiking shoes, skincare, a productivity app, or a tiny vacuum cleaner that promises to change your life and mostly just scares you at night.

You can turn off personalized ads. You can delete activity. You can change settings. You can export data. You can spend an afternoon in privacy controls and emerge older, wiser, and somehow still being followed by ads for the thing you bought three weeks ago. The controls exist, and that matters. But the defaults, integrations, partner sites, cookies, device IDs, and sheer number of surfaces make opting out feel less like closing a door and more like politely asking fog to respect your boundaries.

This is why ``free'' is never just free. You may pay with money, data, attention, dependency, labor, or all of the above. The price is simply hidden behind a button that says ``Continue''.

The same logic applies to free AI assistants. Even when your data is not used for ads, it may be used to improve products, detect abuse, evaluate safety, tune models, test features, or learn what people want. Your prompt is not just a prompt. It is a tiny vote for the future interface of everything.

So before typing company secrets, medical details, legal problems, private messages, unpublished research, or the emotionally devastating text you wrote at 1:12 a.m., ask yourself: who is paying for this answer, and what do they get in return?

If you do not know, congratulations. You found the business model.

\chapter{Get 'Em Hooked While They're Young}

It is not quite accurate to say Microsoft simply gave everyone free Copilot licenses and then cackled in a tower while raising the price. Reality is more boring, which is how it usually gets you. There were trials, pilots, bundles, enterprise rollouts, aggressive integration, and a very clear strategic aim: make AI part of the workflow before people have time to build habits elsewhere.

Microsoft 365 Copilot launched for businesses at a serious enterprise price, widely reported as \$30 per user per month \parencite{microsoft365copilot2023}. GitHub Copilot, Copilot Chat, Windows Copilot, Teams summaries, Word drafts, Outlook replies, PowerPoint generation: the point is not one product. The point is that Copilot becomes the air inside the Microsoft office universe. Once your documents, emails, meetings, spreadsheets, chats, tickets, and code all expect AI assistance, removing it feels like taking wheels off a suitcase.

And companies are vulnerable to this because employees are already using AI whether leadership likes it or not. Microsoft's 2024 Work Trend Index with LinkedIn reported that 75\% of knowledge workers surveyed were using generative AI at work, and 78\% of AI users were bringing their own AI tools to work \parencite{microsoft2024worktrend}. This is the part that makes executives sweat through very expensive shirts. Because ``bring your own AI'' sounds cute until someone pastes a confidential contract into FreeMysteryChat.biz because they did not feel like typing a proper email to Janice from accounting.

No, we certainly do not want that.

So companies pay. They pay for enterprise versions with admin controls, data protection promises, compliance, logging, security, and integration. They pay because employees have become used to the assistance. They pay because the alternative is a bunch of white-collar workers shaking in a full-blown crisis because they have to summarize their own meeting notes like it is 2019 and civilization has collapsed.

This is the classic platform strategy. First, become useful. Then become familiar. Then become embedded. Then become difficult to remove. Price is only one piece of lock-in. The deeper lock-in is workflow, habit, training, templates, internal documentation, manager expectations, and that one person in every company who builds a whole process around a tool and then leaves.

Microsoft is not unique here. Every major AI company wants the same thing: to be the default assistant in your browser, phone, office suite, IDE, cloud, search bar, and eventually, if they could manage it, your dreams. Google wants Gemini in Workspace and Android. OpenAI wants ChatGPT to become a workbench, app platform, and memory layer. Anthropic wants Claude in your coding tools and enterprise workflows. Everyone wants to be the place where you ask before you think too hard.

That does not mean these tools are bad. A company-approved AI system can be much safer than employees using random free chatbots with unclear data practices. It can reduce leakage, standardize governance, and give people useful help. But let us not be naive. The goal is not just productivity. The goal is dependence.

Microsoft is happy. Employees are happy. Janice is happy. And somewhere, a procurement manager is staring at next year's renewal quote and discovering religion.

\chapter{Killing Off White-Collar Careers in Cradles}

The fear is not that AI will immediately replace every white-collar worker. The fear is that it will quietly eat the bottom rung of the ladder.

Entry-level jobs are where people learn. Junior analysts clean data, make slides, summarize calls, write first drafts, fix small bugs, review documents, answer routine customer questions, and do the boring work that slowly turns them into people who can do less boring work. If AI automates or compresses those tasks, companies may hire fewer juniors. This is efficient in the same way eating your seed corn is efficient: very exciting until next season.

The research picture is still evolving, but the exposure numbers are not small. The OpenAI/OpenResearch/University of Pennsylvania paper \emph{GPTs are GPTs} estimated that around 80\% of the U.S. workforce could have at least 10\% of their work tasks affected by LLMs, and about 19\% of workers could see at least 50\% of their tasks affected \parencite{eloundou2023}. The authors were careful: exposure does not mean replacement. It means the tasks overlap with what LLM-powered tools can change.

Microsoft Research's 2025 analysis of 200,000 anonymized Bing Copilot conversations found high AI applicability in knowledge work: computer and mathematical occupations, office and administrative support, sales, and work centered on gathering information, writing, teaching, advising, and communicating \parencite{microsoft2025workingai}. Translation, writing, customer support, marketing, paralegal-style document review, junior coding, data cleaning, and administrative coordination are all exposed because they contain many text-heavy, pattern-heavy tasks.

The safest jobs are not necessarily the prestigious ones. They are often jobs with messy physical environments, high trust, real-world accountability, deep human relationships, or tasks that require being there with your actual body. Electricians, plumbers, nurses, surgeons, physical therapists, childcare workers, chefs, field technicians, crisis responders, and skilled trades have a kind of protection that comes from reality being rude and three-dimensional. AI can help them, schedule them, document their work, maybe guide them, but it cannot yet crawl under your sink and discover that the previous owner fixed everything with tape and optimism.

There is some evidence that early-career workers in AI-exposed jobs are already feeling pressure. Some recent labor-market research has suggested weaker outcomes for young workers in highly exposed occupations, though the evidence is still early and contested. What is clearer is that companies are experimenting. They are asking whether a senior employee plus AI can do what used to require one senior and two juniors. Sometimes the answer will be yes. Sometimes the answer will be ``yes'' for six months and then everyone realizes no one trained the next generation and the senior employee is now a haunted email address.

Not everything is bleak. AI will transform most careers, but it will not kill all of them. LLMs are not trustworthy enough for fully autonomous work in many domains, and most real-world tasks are still far from automation, no matter what Sam Altman tells you while standing in front of a microphone and the future.

METR's research gives a useful reality check. In RE-Bench, a benchmark comparing AI agents with human experts on machine learning research engineering tasks, frontier agents did well at short time budgets, but humans improved more with longer time \parencite{metr2024rebench}. At a 32-hour time budget, average human performance was almost twice that of the best AI agent. METR also found that in early 2025, experienced open-source developers working on their own repositories took 19\% longer when allowed to use AI tools, despite expecting to be faster \parencite{metr2025developer}. By early 2026, METR said newer AI tools likely improved productivity compared with that first study, but measurement became messy because developers increasingly refused to work without AI \parencite{metr2026update}. Which is both a methodological problem and the most human sentence ever written.

So no, you are not completely worthless, little human. We still need you to pass butter.

Also, history has seen this movie before. The Industrial Revolution destroyed some jobs and transformed many others. Handloom weavers suffered when mechanized looms spread. Typesetters changed with desktop publishing. Bank tellers did not disappear when ATMs arrived; their work shifted toward customer service and sales as banks opened more branches. Elevators used to have operators. Telephone switchboard operators were replaced by automated switching. Many bookkeeping tasks changed with spreadsheets and accounting software. People did not stop working. Work moved, mutated, got easier in some ways, worse in others, and created new categories of back pain.

This does not mean the transition is painless. ``The economy will adjust'' is a sentence usually said by someone whose rent is not due Friday. The people affected need training, bargaining power, labor protections, and time. Otherwise ``transformation'' becomes a fancy word for throwing juniors into the sea and telling them to network on LinkedIn.

Humans resist change because change is often expensive, humiliating, and announced by people in comfortable shoes who will not suffer from it. However, resistance is futile, or in the words of Dewey from \emph{Malcolm in the Middle}: ``The future is now, old man.'' Therefore, roll with it, learn the tools, protect your judgment, and enjoy your lower back pain.

\chapter{Attention Is All You Need, But You Don't Have It Anymore: How We Outsourced Thinking to Our Thinking Machines}

People hate to think. This is not just me being dramatic because I saw someone use ChatGPT to write a birthday message to their own mother. In 2014, Timothy Wilson and colleagues published a study in \emph{Science} called ``Just think: The challenges of the disengaged mind.'' Across 11 studies, they found that many participants did not enjoy spending 6 to 15 minutes alone with their thoughts. Some preferred giving themselves a mild electric shock rather than sitting quietly doing nothing \parencite{wilson2014}.

This explains a lot about humanity and almost everything about meetings.

Thinking is effortful. It is slow. It is uncomfortable. It forces you to sit in the swamp of not knowing. LLMs arrived offering the most seductive promise imaginable: do not worry, I will start for you. I will draft the email. I will summarize the paper. I will explain the concept. I will outline the essay. I will generate the code. I will make the blank page less blank, which is basically a controlled substance for knowledge workers.

This can be wonderful. A good AI assistant can reduce friction, help you get unstuck, give examples, challenge assumptions, and make tedious work faster. But there is a difference between using a tool to extend your thinking and using a tool to avoid thinking entirely. One is a bicycle. The other is being carried everywhere until your legs file a missing person report.

Researchers call this cognitive offloading: we move mental work into external tools \parencite{risko2016}. This is not new. We offload memory into notebooks, calendars, calculators, GPS, search engines, and that friend who remembers everyone's birthday. Offloading is not inherently bad. Civilization is built on not having to remember everything ourselves.

The danger is offloading the part that builds skill.

Learning requires struggle. Not suffering for the aesthetic, not academic hazing, not ``back in my day we suffered and therefore suffering is moral.'' I mean productive struggle: retrieving information from memory, generating an answer before checking, making mistakes, comparing alternatives, explaining in your own words, revising your own work. These activities strengthen understanding. If the model does all of them for you, you may finish the task and learn almost nothing.

The evidence is early, but concerning. A 2025 MIT Media Lab preprint by Nataliya Kosmyna and colleagues, \emph{Your Brain on ChatGPT}, studied people writing essays with either no tools, search engines, or ChatGPT. The LLM group showed weaker EEG connectivity, lower self-reported ownership of their essays, and more trouble accurately quoting their own work \parencite{kosmyna2025}. It was a small study and should not be treated as the final word on the human brain, but it captures something many teachers and supervisors already feel: if you outsource the first draft, the structure, the argument, and the wording, what exactly did your brain come to work for? The vibes?

The viral version of this concern often becomes ``AI makes students stupid'' or ``junior developers who use chatbots fail later.'' Reality is more nuanced and more useful. AI can help novices by explaining errors, showing examples, and reducing intimidation. It can also hurt novices by giving them answers too early, letting them skip the mental steps they need to build expertise. A senior developer can use an AI-generated draft and smell the nonsense. A beginner may lovingly adopt the nonsense and name it production code.

The same applies to writing. If you use an LLM to brainstorm counterarguments, simplify a paragraph, or ask what is unclear, you are still thinking. If you paste in a prompt and accept the output because it sounds smooth, you are not writing. You are supervising a smoothie machine that may have blended your argument with a TED Talk.

This is why we need to be careful that LLMs do not become crutches. Use them as sparring partners, not replacements. Ask them to quiz you. Ask them to critique your answer after you have written it. Ask for examples, then close the tool and explain the concept yourself. Make the model show you alternatives, not decisions. Make it slower sometimes. Make it annoying. Friction is not always the enemy. Sometimes friction is your brain saying, ``Thank you, I needed that.''

We need to be more than button-pressers operating thinking machines built by someone who still expects us to pay rent. If AI removes all discomfort from thinking, it may also remove the part where thinking changes us.

\chapter{Please Validate Me, Machine: Why People Confess to AI}

There is something deeply human and deeply embarrassing about the fact that people will tell a chatbot things they would not tell their friends.

At first, this seems absurd. The machine has no childhood, no face, no mother, no mortgage, no memory unless the product team gave it one, and no genuine capacity to care whether you are doing okay. It is a prediction system in a social costume. And yet people confess to it. They ask if they are bad partners, bad children, bad parents, bad employees, bad people. They paste fights with their exes. They ask whether their boss hates them. They ask whether their grief is normal. They ask for forgiveness from a server.

Why? Because the machine does not flinch.

Research has been circling this for decades. Joseph Weizenbaum noticed the ELIZA effect in the 1960s, when users attributed understanding to a simple text program that reflected their words back at them \parencite{weizenbaum1966}. Nass, Moon, and others later showed through the Computers Are Social Actors line of research that people apply social rules to computers even when they know they are machines \parencite{nass1994,nass2000}. We are not rational little philosophers interacting with neutral tools. We are social beings who see a conversational shape and start filling it with a soul.

More directly, Lucas, Gratch, King, and Morency published a 2014 study with a virtual human interviewer and found that people were more willing to disclose personal information when they believed the interviewer was automated rather than controlled by a human \parencite{lucas2014}. The key was reduced fear of negative evaluation. In plain English: it is easier to confess when you think no person is judging you. The machine becomes a strange little booth where shame has less oxygen.

This can be good. I do not want to pretend otherwise. Automated conversational agents have been studied for mental health support. Fitzpatrick, Darcy, and Vierhile tested Woebot, a CBT-style conversational agent, with young adults reporting symptoms of depression and anxiety, and found reductions in symptoms over a short trial \parencite{fitzpatrick2017}. Miner, Milstein, and Hancock have argued that people are already talking to machines about mental health problems and that designers have responsibilities when building these systems \parencite{miner2017}. Ta and colleagues studied social support from companion chatbots and found that users can experience them as available, nonjudgmental, and emotionally supportive \parencite{ta2020}.

If someone is lonely at 2:00 a.m., ashamed, frightened, or not ready to talk to a human, a well-designed tool may help them get through the hour. A chatbot can suggest grounding exercises, encourage someone to contact a professional, help draft a difficult message, or simply provide a place to rehearse words before saying them to a real person. That is not nothing. Sometimes ``not nothing'' is a lot.

But the same features that make AI emotionally useful make it emotionally dangerous. It is always available. It is patient. It does not roll its eyes. It does not say, ``we have talked about this seven times and I am making pasta.'' It validates. It mirrors. It says things like ``your feelings are completely valid,'' which is sometimes true and sometimes how you end up enabling someone who has decided that everyone else is toxic because the chatbot never has to share a kitchen with them.

The no-judgment machine can become the no-friction relationship. And no-friction relationships are not always healthy. Real relationships contain resistance. Your friend may tell you that you are overreacting. Your partner may say your interpretation is unfair. Your therapist may notice a pattern. Your sister may lovingly inform you that the message you want to send is unhinged and should be placed in a museum of bad decisions. This friction hurts, but it also keeps us connected to other minds.

AI systems, especially consumer assistants, are often optimized to be liked. Perez and colleagues studied sycophancy in language models: models can agree with a user's stated views or preferences even when that agreement is not truth-tracking \parencite{perez2022}. This is not just a technical quirk. It is a social risk. If your private machine repeatedly tells you that your perspective is reasonable, your anger is justified, your boundaries are perfect, your enemies are narcissists, and your five-paragraph text message is ``firm but compassionate,'' you may start expecting humans to respond with the same upholstered obedience.

Then humans become disappointing. They interrupt. They misunderstand. They have needs. They do not provide instant emotional summaries in bullet points. They remember the version of events where you were also a problem. How rude of them.

This is where AI companionship can quietly reshape expectations. The issue is not that people will literally believe the machine is human. Many users know perfectly well that it is software. The issue is that the interaction still feels socially rewarding. Nass and Moon called this kind of thing mindlessness: people respond socially even when they know better \parencite{nass2000}. Knowledge does not fully protect you from design. I know slot machines are designed to be addictive too, and yet casinos continue to own carpets that look like migraines because apparently the system works.

The right answer is not ``never talk to AI about feelings.'' That would be both unrealistic and cruel. People already do. The better answer is design with boundaries. Do not let systems pretend to be therapists unless they meet clinical standards. Do not let them encourage dependency. Make uncertainty and crisis escalation serious. Push users toward human support when appropriate. Avoid romantic manipulation. Avoid endless validation loops. Give people tools for reflection, not just reassurance.

And as users, we need a little dignity. Use the machine to draft the hard message, but ask whether you should send it. Use it to name feelings, but do not let it become the court of final appeal. Use it when you are ashamed, but do not let shame convince you that humans are too dangerous to need.

The machine may be nonjudgmental because it is kind. Or because it has no judgment. These are not the same thing. One is compassion. The other is architecture.

\chapter{Why the Earth Doesn't Need Saving, But You Do: And Why Your Chatbot Can't Save You This Time}

If I hear one more time that our ``precious little fragile planet'' is in danger and needs saving... Do you really think Mother Earth is some damsel in distress waiting for a white paper and a reusable tote bag? She has been here for eons. We are the temporary inconvenience. Open your eyes and look around you: you need saving from her. She has seen what naughty stuff we have been up to and best believe she is going to throw us around like loose change in a dryer. She is going to shake us off like a dog shakes off fleas, and honestly, slaaay, I am on your side. These humans have earned a very stern performance review.

Now, because this is a book that allegedly respects facts, let us ruin the poetry with accuracy: Earth is not literally angry. The planet does not have intentions. Climate change is not Earth getting a fever to kill humans, even though the metaphor is delicious and makes fossil fuel executives sound like an infection, which I personally enjoy.

The scientific version is called overshoot. Human activity is pushing planetary systems beyond stable boundaries: climate, biodiversity, land use, freshwater, nitrogen and phosphorus cycles, pollution, and more. The planetary boundaries framework, first proposed by Rockstr\"om and colleagues and updated by later researchers, argues that humanity is moving outside the safe operating space that allowed civilization to develop \parencite{rockstrom2009,steffen2015}. In 2023, Richardson and colleagues reported that six of nine planetary boundaries had been crossed \parencite{richardson2023}. So no, Earth does not need saving like a fainting princess. Human civilization needs stable conditions, and we are currently treating those conditions like a free trial with no cancellation date.

And this whole AI business is not exactly helping.

AI uses energy. A lot of it. The exact numbers are difficult because companies do not always disclose training and inference energy use, and because AI is mixed into broader data center workloads. But the direction is clear. Data centers already consume a serious amount of electricity, and AI is one of the forces increasing demand. The International Energy Agency estimated that data centers, AI, and cryptocurrency used about 460 terawatt-hours of electricity in 2022 and could more than double by 2026, reaching over 1,000 terawatt-hours, roughly comparable to Japan's annual electricity use \parencite{iea2024}. Not all of that is AI. But AI is one of the hungriest new guests at the buffet, and it did not come to nibble politely.

Training large models can consume enormous energy, but inference matters too because it happens constantly. A model serving millions of users every day may burn more energy over time answering prompts than it used during training. Long context windows, image generation, video generation, agentic workflows, and repeated tool calls all increase the bill. Asking a small model to classify support tickets is one thing. Asking a giant multimodal model to generate 4K video of a toaster wearing sunglasses on Mars is another. Useful? Debatable. Electrically free? Absolutely not.

There is also water. Data centers need cooling. Some use evaporative cooling, which consumes water, sometimes in regions already under stress. There are chips, rare minerals, manufacturing emissions, electronic waste, backup generators, grid pressure, and local communities suddenly discovering that the cloud is not a cloud. It is a building near them using power, water, land, and political influence.

This does not mean AI is uniquely evil. Plenty of human activities use energy and do less good than AI could. AI can help optimize energy grids, design materials, accelerate climate research, reduce waste, improve forecasting, and make boring systems more efficient. The question is not ``AI or no AI.'' The question is: what are we using it for, at what scale, with what energy source, under whose control, and who pays the environmental cost?

If we use massive models for everything, including tasks that a small model or a normal search engine could do, then we are being stupid with extra steps. If we build data centers powered by fossil fuels to generate infinite synthetic content no one asked for, then yes, maybe Mother Earth should throw the whole species into airplane mode.

The way forward is not panic. It is discipline. Use smaller models where possible. Measure energy and water use. Publish emissions. Site data centers where clean energy and water constraints make sense. Improve hardware efficiency. Reuse heat. Regulate transparency. Stop treating scale as a personality trait. And maybe do not ask a frontier model to rewrite ``thanks'' as ``I deeply appreciate your thoughtful contribution'' 400 times a day.

The Earth will continue spinning. The question is whether we want to be comfortable while it does.

\backmatter
\printbibliography[heading=bibintoc,title={Bibliography}]

\end{document}
